\section{R ile test gücü ve tek örneklem testi}
\label{sec:r-b10}

\textbf{Soru.} $n=25$, $\bar x=55$, $s=10$ için referans ortalama 50'yi
çift yönlü test edin. Ayrıca veri toplanmadan önce gerçek farkın 5, evren
standart sapmasının 10 olduğu varsayımıyla $n=34$ için gücü hesaplayın.

\begin{lstlisting}[style=prismR]
hacim <- 25
ortalama <- 55
standart_sapma <- 10
referans <- 50
standart_hata <- standart_sapma / sqrt(hacim)
t_degeri <- (ortalama - referans) / standart_hata
p_degeri <- 2 * pt(abs(t_degeri), df = hacim - 1,
                   lower.tail = FALSE)
aralik <- ortalama + c(-1, 1) *
  qt(0.975, df = hacim - 1) * standart_hata
c(t = t_degeri, p = p_degeri,
  d = (ortalama - referans) / standart_sapma)
aralik
power.t.test(n = 34, delta = 5, sd = 10,
             sig.level = 0.05, type = "one.sample",
             alternative = "two.sided", strict = TRUE)
plan <- power.t.test(delta = 5, sd = 10, power = 0.80,
                     sig.level = 0.05, type = "one.sample",
                     alternative = "two.sided", strict = TRUE)
ceiling(plan$n)
\end{lstlisting}

\begin{outputguidebox}
\textbf{Çözüm ve kontrol değerleri.} $t(24)=2{,}50$, $p=0{,}01965$,
$d=0{,}50$ ve yüzde 95 güven aralığı $[50{,}8722;59{,}1278]$'dir.
Planlama varsayımları altında 34 gözlemin gücü yaklaşık 0,8078'dir.
Yüzde 80 hedef güç için hesaplanan hacim yukarı yuvarlandığında 34 bulunur.
\end{outputguidebox}

\textbf{Yorum.} \texttt{strict = TRUE} çift yönlü testin iki ret bölgesini
de güç hesabına katar. Güç hesabındaki 5 puan ve 10 puan önceden gerekçelenmiş
planlama değerleri olmalıdır; gözlenen anlamlılığı yeniden ifade eden
``gözlenen güç'' olarak sunulmamalıdır.

\textbf{Pekiştirme ve yanıt.} Tip II hata olasılığı bu belirli alternatif
ve tasarım altında yaklaşık $1-0{,}8078=0{,}1922$'dir. Daha küçük bir gerçek
fark için aynı örneklem hacminin gücü daha düşük olur.